% Part: first-order-logic
% Chapter: natural-deduction
% Section: proving-things

\documentclass[../../../include/open-logic-section]{subfiles}

\begin{document}

\iftag{FOL}
      {\olfileid{fol}{ntd}{pro}}
      {\olfileid{pl}{ntd}{pro}}

\olsection{\usetoken{P}{derivation}: उदाहरणे}

\begin{ex}
!!{sentence} $(!A \land !B) \lif !A$ ची !!a{derivation} देऊ.

हवा असलेला निष्कर्ष !!{derivation} च्या तळाशी लिहून आपण सुरुवात करतो.
\begin{prooftree}
\AxiomC{}
\UnaryInfC{$(!A\land !B) \lif !A$}
\end{prooftree}

पुढे, या रूपाचे !!a{sentence} कोणत्या प्रकारच्या अनुमानामुळे मिळू शकते हे
शोधावे लागते. निष्कर्षाचा !!{main operator} हा $\lif$ आहे; म्हणून
\Intro{\lif} नियम वापरून निष्कर्षापर्यंत पोहोचण्याचा प्रयत्न करू. काम पुढे
जाताना संबंधित गृहीतके लिहून ठेवणे आणि निगमन नियमांना लेबले देणे उत्तम;
त्यामुळे सिद्धतेच्या शेवटी सर्व गृहीतके !!{discharged} झाली आहेत का हे सहज
दिसते.
\begin{prooftree}
\AxiomC{$\Discharge{!A \land !B}{1}$}
\DeduceC{$!A$}
\DischargeRule{\Intro{\lif}}{1}
\UnaryInfC{$(!A\land !B) \lif !A$}
\end{prooftree}

आता गृहीतक $!A \land !B$ पासून $!A$ पर्यंतच्या पायऱ्या भराव्या लागतील.
येथे हाताळण्यासाठी फक्त एकच संयोजक, $\land$, असल्यामुळे $\land$ चा विलोपन
नियम वापरलाच पाहिजे. त्यामुळे पुढील सिद्धता मिळते:
\begin{prooftree}
\AxiomC{$\Discharge{!A \land !B}{1}$}
\RightLabel{\Elim{\land}}
\UnaryInfC{$!A$}
\DischargeRule{\Intro{\lif}}{1}
\UnaryInfC{$(!A\land !B) \lif !A$}
\end{prooftree}
आता आपल्याकडे $(!A \land !B) \lif !A$ ची योग्य !!{derivation} आहे.
\end{ex}

\begin{ex}
आता $(\lnot !A \lor !B) \lif (!A \lif !B)$ ची !!a{derivation} देऊ.

हवा असलेला निष्कर्ष !!{derivation} च्या तळाशी लिहून आपण सुरुवात करतो.
\begin{prooftree}
\AxiomC{}
\UnaryInfC{$(\lnot !A \lor !B) \lif (!A \lif !B)$}
\end{prooftree}
हा निष्कर्ष देऊ शकणारा तार्किक नियम शोधण्यासाठी निष्कर्षातील $\lnot$,
$\lor$ आणि $\lif$ हे तार्किक संयोजक पाहतो. सध्या $\lif$ ची पहिली आस्थितीच
महत्त्वाची आहे, कारण अंतिम क्रमवर्तीतील !!{sentence} चा तो !!{main operator}
आहे; $\lnot$, $\lor$ आणि $\lif$ ची दुसरी आस्थिती दुसऱ्या संयोजकाच्या
व्याप्तीत असल्यामुळे त्यांचा विचार नंतर करू. म्हणून आपण \Intro{\lif}
नियमाने सुरुवात करतो. त्याचे योग्य उपयोजन पुढीलप्रमाणे असले पाहिजे:
\begin{prooftree}
\AxiomC{$\Discharge{\lnot !A \lor !B}{1}$}
\DeduceC{$!A \lif !B$}
\DischargeRule{\Intro{\lif}}{1}
\UnaryInfC{$(\lnot !A \lor !B) \lif (!A \lif !B)$}
\end{prooftree}
पुढे जाण्यासाठी आता दोन पर्याय आहेत. एकतर तळापासून वरच्या दिशेने काम
चालू ठेवून \Intro{\lif} नियमाचे आणखी एक उपयोजन शोधता येईल, किंवा वरून
खाली काम करून \Elim{\lor} नियम लागू करता येईल. दुसरा पर्याय घेऊ.
\Elim{\lor} चे सर्वांत डावे आधारविधान म्हणून $\lnot !A \lor !B$ हे
गृहीतक वापरू. \Elim{\lor} च्या योग्य उपयोजनासाठी इतर दोन्ही आधारविधाने
$!A \lif !B$ या निष्कर्षाशी एकरूप असली पाहिजेत; मात्र प्रत्येकाला
अनुक्रमे आणखी एका गृहीतकापासून, म्हणजे $\lnot !A \lor !B$ च्या दोन
विकल्पघटकांपैकी एकापासून, निष्पन्न करता येते. त्यामुळे आपली
!!{derivation} पुढीलप्रमाणे दिसेल:
\begin{prooftree}
\AxiomC{$\Discharge{\lnot !A \lor !B}{1}$}
\AxiomC{$\Discharge{\lnot !A}{2}$}
\DeduceC{$!A \lif !B$}
\AxiomC{$\Discharge{!B}{2}$}
\DeduceC{$!A \lif !B$}
\DischargeRule{\Elim{\lor}}{2}
\TrinaryInfC{$!A \lif !B$}
\DischargeRule{\Intro{\lif}}{1}
\UnaryInfC{$(\lnot !A \lor !B) \lif (!A \lif !B)$}
\end{prooftree}

उजवीकडील दोन्ही शाखांत आपल्याला $!A \lif !B$ हे !!{derive} करायचे आहे;
त्यासाठी \Intro{\lif} वापरणे सर्वांत योग्य आहे.
\begin{prooftree}
\AxiomC{$\Discharge{\lnot !A \lor !B}{1}$}
\AxiomC{$\Discharge{\lnot !A}{2}, \Discharge{!A}{3}$}
\DeduceC{$!B$}
\DischargeRule{\Intro{\lif}}{3}
\UnaryInfC{$!A \lif !B$}
\AxiomC{$\Discharge{!B}{2}, \Discharge{!A}{4}$}
\DeduceC{$!B$}
\DischargeRule{\Intro{\lif}}{4}
\UnaryInfC{$!A \lif !B$}
\DischargeRule{\Elim{\lor}}{2}
\TrinaryInfC{$!A \lif !B$}
\DischargeRule{\Intro{\lif}}{1}
\UnaryInfC{$(\lnot !A \lor !B) \lif (!A \lif !B)$}
\end{prooftree}

!!{derivation} च्या उरलेल्या दोन भागांसाठी, मधल्या भागात $\lnot !A$ आणि
$!A$ यांपासून, तर उजवीकडील भागात $!A$ आणि $!B$ यांपासून, $!B$ मिळवणाऱ्या
!!{derivation}s हवी आहेत. यांपैकी पहिली आधी घेऊ. $\lnot !A$ आणि $!A$ ही
\Elim{\lnot} ची दोन आधारविधाने आहेत:
\begin{prooftree}
\AxiomC{$\Discharge{\lnot !A}{2}$}
\AxiomC{$\Discharge{!A}{3}$}
\RightLabel{\Elim{\lnot}}
\BinaryInfC{$\lfalse$}
\DeduceC{$!B$}
\end{prooftree}
\FalseInt वापरून $!B$ हा निष्कर्ष मिळवता येतो आणि शाखा पूर्ण करता येते.
\begin{prooftree}
\AxiomC{$\Discharge{\lnot !A \lor !B}{1}$}
\AxiomC{$\Discharge{\lnot !A}{2}$}
\AxiomC{$\Discharge{!A}{3}$}
\RightLabel{\Intro{\lfalse}}
\BinaryInfC{$\lfalse$}
\RightLabel{\FalseInt}
\UnaryInfC{$!B$}
\DischargeRule{\Intro{\lif}}{3}
\UnaryInfC{$!A \lif !B$}
\AxiomC{$\Discharge{!B}{2}, \Discharge{!A}{4}$}
\DeduceC{$!B$}
\DischargeRule{\Intro{\lif}}{4}
\UnaryInfC{$!A \lif !B$}
\DischargeRule{\Elim{\lor}}{2}
\TrinaryInfC{$!A \lif !B$}
\DischargeRule{\Intro{\lif}}{1}
\UnaryInfC{$(\lnot !A \lor !B) \lif (!A \lif !B)$}
\end{prooftree}

आता सर्वांत उजवी शाखा पाहू. येथे हे लक्षात घेणे महत्त्वाचे आहे की
!!{derivation} ची व्याख्या
\emph{गृहीतके मुक्त करण्याची परवानगी देते}, पण
\emph{तसे करणे आवश्यक ठरवत नाही}. दुसऱ्या शब्दांत, $!A$ आणि $!B$ या
गृहीतकांपैकी एकाचा उपयोग न करता दुसऱ्यापासून $!B$ निष्पन्न करता येत असेल,
तर ते योग्य आहे. $!B$ पासून $!B$ हे !!{derive} करणे तर क्षुल्लक आहे:
एकटे $!B$ ही अशी !!a{derivation} असून कोणत्याही अनुमानाची गरज नाही. म्हणून
गृहीतक~$!A$ सरळ वगळता येते.
\begin{prooftree}
\AxiomC{$\Discharge{\lnot !A \lor !B}{1}$}
\AxiomC{$\Discharge{\lnot !A}{2}$}
\AxiomC{$\Discharge{!A}{3}$}
\RightLabel{\Elim{\lnot}}
\BinaryInfC{$\lfalse$}
\RightLabel{\FalseInt}
\UnaryInfC{$!B$}
\DischargeRule{\Intro{\lif}}{3}
\UnaryInfC{$!A \lif !B$}
\AxiomC{$\Discharge{!B}{2}$}
\RightLabel{\Intro{\lif}}
\UnaryInfC{$!A \lif !B$}
\DischargeRule{\Elim{\lor}}{2}
\TrinaryInfC{$!A \lif !B$}
\DischargeRule{\Intro{\lif}}{1}
\UnaryInfC{$(\lnot !A \lor !B) \lif (!A \lif !B)$}
\end{prooftree}
लक्षात घ्या की पूर्ण झालेल्या !!{derivation} मध्ये सर्वांत उजवीकडील
\Intro{\lif} अनुमान प्रत्यक्षात कोणतेही गृहीतक मुक्त करत नाही.
\end{ex}

\begin{ex}
आतापर्यंत आपल्याला \FalseCl{} नियमाची गरज पडलेली नाही. हा नियम विशेष आहे,
कारण नियमाच्या निष्कर्षाचे उप-!!{formula} नसलेले गृहीतक तो मुक्त करू देतो.
त्याचा \FalseInt{} नियमाशी जवळचा संबंध आहे. खरे तर \FalseInt{} नियम हे
\FalseCl{} नियमाचे विशेष प्रकरण आहे---``अंतःप्रज्ञावादी तर्कशास्त्र''
नावाचे एक तर्कशास्त्र आहे आणि त्यात केवळ \FalseInt{} अनुज्ञात आहे. इतर
कोणताही मार्ग चालत नसताना \FalseCl{} नियम हा अखेरचा उपाय असतो. उदाहरणार्थ,
आपल्याला $!A \lor \lnot !A$ हे !!{derive} करायचे आहे असे समजा. नेहमीच्या
पद्धतीनुसार $\Intro{\lor}$ वापरून $!A \lor \lnot !A$ हे !!{derive}
करण्याचा प्रयत्न करू. पण त्यासाठी कोणत्याही गृहीतकांशिवाय $!A$ किंवा
$\lnot !A$ यांपैकी एक !!{derive} करावे लागेल, आणि ते शक्य नाही. अशा वेळी
\FalseCl{} उपयोगी पडतो!
\begin{prooftree}
  \AxiomC{$\Discharge{\lnot(!A \lor \lnot !A)}{1}$}
  \DeduceC{$\lfalse$}
  \DischargeRule{\FalseCl}{1}
  \UnaryInfC{$!A \lor \lnot !A$}
\end{prooftree}
आता $\lnot(!A \lor \lnot !A)$ पासून $\lfalse$ ची !!a{derivation} शोधायची
आहे. $\lfalse$ हा $\Elim{\lnot}$ चा निष्कर्ष असल्यामुळे तो नियम वापरून
पाहू:
\begin{prooftree}
  \AxiomC{$\Discharge{\lnot(!A \lor \lnot !A)}{1}$}
  \DeduceC{$\lnot !A$}
  \AxiomC{$\Discharge{\lnot(!A \lor \lnot !A)}{1}$}
  \DeduceC{$!A$}
  \RightLabel{\Elim{\lnot}}
  \BinaryInfC{$\lfalse$}
  \DischargeRule{\FalseCl}{1}
  \UnaryInfC{$!A \lor \lnot !A$}
\end{prooftree}
$\lnot !A$ ची !!a{derivation} शोधण्याच्या आपल्या पद्धतीनुसार
$\Intro{\lnot}$ चे उपयोजन करावे लागेल:
\begin{prooftree}
  \AxiomC{$\Discharge{\lnot(!A \lor \lnot !A)}{1}, \Discharge{!A}{2}$}
  \DeduceC{$\lfalse$}
  \DischargeRule{\Intro{\lnot}}{2}
  \UnaryInfC{$\lnot !A$}
  \AxiomC{$\Discharge{\lnot(!A \lor \lnot !A)}{1}$}
  \DeduceC{$!A$}
  \RightLabel{\Elim{\lnot}}
  \BinaryInfC{$\lfalse$}
  \DischargeRule{\FalseCl}{1}
  \UnaryInfC{$!A \lor \lnot !A$}
\end{prooftree}
येथे $\lnot(!A \lor \lnot !A)$ हे गृहीतक आणि आपल्या नव्या $!A$
गृहीतकापासून~$\Intro{\lor}$ ने निष्पन्न होणारे $!A \lor \lnot !A$
यांना $\Elim{\lnot}$ लावून $\lfalse$ सहज मिळवता येते:
\begin{prooftree}
  \AxiomC{$\Discharge{\lnot(!A \lor \lnot !A)}{1}$}
  \AxiomC{$\Discharge{!A}{2}$}
  \RightLabel{\Intro{\lor}}
  \UnaryInfC{$!A \lor \lnot !A$}
  \RightLabel{\Elim{\lnot}}
  \BinaryInfC{$\lfalse$}
  \DischargeRule{\Intro{\lnot}}{2}
  \UnaryInfC{$\lnot !A$}
  \AxiomC{$\Discharge{\lnot(!A \lor \lnot !A)}{1}$}
  \DeduceC{$!A$}
  \RightLabel{\Elim{\lnot}}
  \BinaryInfC{$\lfalse$}
  \DischargeRule{\FalseCl}{1}
  \UnaryInfC{$!A \lor \lnot !A$}
\end{prooftree}
उजव्या बाजूला हीच पद्धत वापरतो; फक्त तेथे $!A$ हे~\FalseCl ने मिळवतो:
\begin{prooftree}
  \AxiomC{$\Discharge{\lnot(!A \lor \lnot !A)}{1}$}
  \AxiomC{$\Discharge{!A}{2}$}
  \RightLabel{\Intro{\lor}}
  \UnaryInfC{$!A \lor \lnot !A$}
  \RightLabel{\Elim{\lnot}}
  \BinaryInfC{$\lfalse$}
  \DischargeRule{\Intro{\lnot}}{2}
  \UnaryInfC{$\lnot !A$}
  \AxiomC{$\Discharge{\lnot(!A \lor \lnot !A)}{1}$}
  \AxiomC{$\Discharge{\lnot !A}{3}$}
  \RightLabel{\Intro{\lor}}
  \UnaryInfC{$!A \lor \lnot !A$}
  \RightLabel{\Elim{\lnot}}
  \BinaryInfC{$\lfalse$}
  \DischargeRule{\FalseCl}{3}
  \UnaryInfC{$!A$}
  \RightLabel{\Elim{\lnot}}
  \BinaryInfC{$\lfalse$}
  \DischargeRule{\FalseCl}{1}
  \UnaryInfC{$!A \lor \lnot !A$}
\end{prooftree}
\end{ex}

\begin{prob}
पुढील गोष्टी दाखवणाऱ्या !!{derivation}s द्या:
\begin{enumerate}
\item $!A \land (!B \land !C) \Proves (!A \land !B) \land !C$.
\item $!A \lor (!B \lor !C) \Proves (!A \lor !B) \lor !C$.
\item $!A \lif (!B \lif !C) \Proves !B \lif (!A \lif !C)$.
\item $!A \Proves \lnot\lnot !A$.
\end{enumerate}
\end{prob}

\begin{prob}
पुढील गोष्टी दाखवणाऱ्या !!{derivation}s द्या:
\begin{enumerate}
\item $(!A \lor !B) \lif !C \Proves !A \lif !C$.
\item $(!A \lif !C) \land (!B \lif !C) \Proves (!A \lor !B) \lif !C$.
\item $\Proves \lnot(!A \land \lnot !A)$.
\item $!B \lif !A \Proves \lnot !A \lif \lnot !B$.
\item $\Proves (!A \lif \lnot !A) \lif \lnot !A$.
\item $\Proves \lnot(!A \lif !B) \lif \lnot !B$.
\item $!A \lif !C \Proves \lnot (!A \land \lnot !C)$.
\item $!A \land \lnot !C \Proves \lnot (!A \lif !C)$.
\item $!A \lor !B, \lnot !B \Proves !A$.
\item $\lnot !A \lor \lnot !B \Proves \lnot(!A \land !B)$.
\item $\Proves (\lnot !A \land \lnot !B) \lif\lnot(!A \lor !B)$.
\item $\Proves \lnot(!A \lor !B) \lif (\lnot !A \land \lnot !B)$.
\end{enumerate}
\end{prob}

\begin{prob}
पुढील गोष्टी दाखवणाऱ्या !!{derivation}s द्या:
\begin{enumerate}
\item $\lnot(!A \lif !B) \Proves !A$.
\item $\lnot(!A \land !B) \Proves \lnot !A \lor \lnot !B$.
\item $!A \lif !B \Proves \lnot !A \lor !B$.
\item $\Proves \lnot \lnot !A \lif !A$.
\item $!A \lif !B, \lnot !A \lif !B \Proves !B$.
\item $(!A \land !B) \lif !C \Proves (!A \lif !C) \lor (!B \lif !C)$.
\item $(!A \lif !B) \lif !A \Proves !A$.
\item $\Proves (!A \lif !B) \lor (!B \lif !C)$.
\end{enumerate}
(या सर्वांसाठी $\FalseCl$~नियम आवश्यक आहे.)
\end{prob}

\end{document}
